% SGA 6, Exposición XIV: fragmento español alineado con la fuente.
% Alcance: secciones 7 y 8, bibliografía y notas finales, hasta la línea
% 15922 inclusive; se detiene antes del índice terminológico.
% Fuente francesa: sga6_fr_r1/src/fr.tex, líneas 15826--15922.
% SHA-256 de la fuente francesa completa:
% 0D2158B9FE3B466309A92165727A0B191422E05FBAA0DACD5806E8322C54BA1A.

\section*{7. Clases de Chern de los complejos perfectos}

\textbf{7.1.} Cuando se plantea la cuestión de desarrollar una fórmula
cohomológica de Riemann--Roch, con o sin denominadores, prescindiendo de la
hipótesis bastante artificial de la existencia de un módulo invertible
amplio sobre los esquemas considerados, uno se enfrenta necesariamente a la
cuestión de definir, para un complejo de módulos perfectos $L^\bullet$ sobre
un esquema $X$, las clases de Chern
\[
 c_i(L_\bullet)\in H^{2i}(X,\Zell(i)),
\]
que generalicen las clases de Chern de los módulos localmente libres de SGA
5 VII. Por otra parte, la misma cuestión se plantea para topoi localmente
anillados generales (compárese [10]), así como para los complejos perfectos
sobre espacios paracompactos anillados por las funciones complejas
continuas, trabajando entonces con la cohomología entera ordinaria. En el
caso de un espacio compacto, se sabe que el $K'(X)$ construido con los
complejos perfectos es el mismo que el «~ingenuo~», construido con los
módulos localmente libres (Exp. IV), y en este caso nos reducimos
trivialmente a la definición clásica. Pero ya en el caso de un espacio
localmente compacto el problema planteado parece no trivial, puesto que la
cohomología entera de tal espacio no siempre es el límite proyectivo de las
cohomologías de sus partes compactas.

\textbf{7.2.} Más precisamente, cuando $Y$ es una parte cerrada de $X$ tal
que $L_\bullet|(X-Y)$ sea nulo (en el sentido de las categorías derivadas),
es decir, tal que los haces de cohomología de $L^\bullet$ tengan soporte en
$Y$, se querrían definir clases de Chern con soporte en $Y$
\[
 c_i^Y(L_\bullet)\in H_Y^{2i}(X,\Zell(i))\quad\text{resp.}\quad H_Y^{2i}(X,\Z),
\]
que restituyeran las clases de Chern «~ordinarias~» (sic) de 7.1 mediante el
homomorfismo canónico
\[
 H_Y^\bullet(X,-)\longrightarrow H^\bullet(X,-).
\]
Incluso en el caso en que estas últimas se reducen a las clases de Chern de
los módulos localmente libres ($X$ esquema con módulo invertible amplio, o
espacio compacto), no se han construido aún estas clases de Chern «~con
soportes~», y sin duda solo podrán construirse al mismo tiempo que las clases
de Chern generales postuladas en 7.1.\footnote{Añadido en octubre de 1968.
Los problemas planteados aquí fueron resueltos en el marco de la cohomología
de Hodge por L. ILLUSIE; cf. [13].}

Es probable que las clases de Chern con soportes tengan un papel que
desempeñar en la demostración de la fórmula cohomológica de Riemann--Roch sin
denominadores (cf. 3.5). Por otra parte, deberían permitir dar una versión
mejorada de la fórmula de Riemann--Roch (con o sin denominadores),
introduciendo en ella familias de soportes.

\section*{8. Anillo de Chow de los esquemas regulares}

El problema consiste en definir la equivalencia racional de los ciclos y
estudiar el anillo de clases obtenido (en particular, sus propiedades
funtoriales), para un esquema noetheriano regular $X$ que admita un módulo
invertible amplio, generalizando la teoría conocida cuando $X$ es liso sobre
un cuerpo [4]. El resultado técnico clave seguirá siendo un «~lema de
desplazamiento~», que asegure que se pueda hacer variar un ciclo dentro de
su clase de modo que corte un ciclo dado con la dimensión correcta.

El problema considerado es bastante natural a priori, tanto más cuanto que
hemos visto (4.5 a 4.8) que los sustitutos $\operatorname{Gr}'(X)$,
$\Gtop^\bullet(X)$ del anillo de Chow $A(X)$ son poco satisfactorios cuando
no nos limitamos a una teoría «~módulo torsión~». Cuando $X$ se obtiene de
un esquema algebraico proyectivo y liso $Y$ sobre un cuerpo $k$, tomando el
esquema local en el vértice de un cono \sisi{proyectante}{proyectivo} $C$ y eliminando el % Canon SGA-SRC-001456
origen, se obtiene (admitiendo que $A(X)$ es efectivamente el límite
inductivo de los anillos de Chow de los $U-(O)$, donde $U$ recorre los
entornos abiertos en $C$ del vértice $O$) que $A(X)$ es canónicamente
isomorfo al anillo cociente $A(Y)/\xi A(Y)$, donde $\xi\in A^1(Y)$ está
inducido por la inmersión proyectiva considerada de $Y$. Esto muestra que el
conocimiento del anillo de Chow $A(X)$ (de tipo no clásico) equivale al
conocimiento de $A(Y)/\xi A(Y)$, que podría denominarse la «~parte
primitiva~» del anillo de Chow $A(Y)$ (en un sentido sugerido por la teoría
cohomológica de Lefschetz--Hodge de las variedades algebraicas proyectivas
no singulares). Este ejemplo muestra que los «~anillos de Chow locales~» de
los anillos locales con singularidades aisladas son invariantes geométricos
nada triviales y que sin duda tendrán un papel que desempeñar en el estudio
«~aritmético~» de estos anillos locales.

Naturalmente, una vez que se disponga de una noción razonable de anillo de
Chow para un esquema noetheriano regular que admita un módulo invertible
amplio, habrá que generalizar a este caso las relaciones establecidas, resp.
conjeturadas, en el n.º 4 entre dicho anillo y el anillo
$\Gtop^\bullet(X)$, y formular en términos de este anillo la cuestión de la
«~fórmula de Riemann--Roch sin denominadores~» enunciada en el n.º 3.

\subsection*{Bibliografía de la Exposición XIV}
\begin{enumerate}[label={[\arabic*]}]
\item M. F. Atiyah y F. Hirzebruch, \emph{The Riemann--Roch theorem for analytic embeddings}, Topology, vol. 1, pp. 151--166 (1962).
\item A. Borel y J.-P. Serre, \emph{Le théorème de Riemann--Roch}, Bull. Soc. Math. France, vol. 86, pp. 97--136 (1958).
\item R. Bott y H. Samelson, \emph{Applications of the theory of Morse to symmetric spaces}, Amer. J. Math., vol. 80, pp. 964--1029 (1958).
\item C. Chevalley, \emph{Les classes d'équivalence rationnelle}, I y II, Séminaire à l'É.N.S., 1958, exposiciones 2 y 3.
\item C. Chevalley, Classification des groupes de Lie algébriques, Séminaire à l'É.N.S., 1956/58.
\item A. Dold y D. Puppe, Homologie nicht-additiver Funktoren. Anwendungen, Annales de l'Institut Fourier, t. XI, pp. 201--312 (1961).
\item H. Grauert, \emph{Ein Theorem der analytischen Garbentheorie}, Publ. Math. IHÉS n.º 5, pp. 5--64 (1960).
\item A. Grothendieck, Classes de faisceaux et théorème de Riemann--Roch, notas multicopiadas del Institute for Advanced Study, 1957, citado aquí como RRR y reproducido en este seminario como apéndice al Capítulo 0.
\item A. Grothendieck, \emph{La théorie des classes de Chern}, Bull. Soc. Math. France, vol. 86, pp. 137--154 (1958).
\item A. Grothendieck, Classes de Chern et représentations linéaires des groupes discrets, en: Dix exposés sur la cohomologie des schémas, Masson--North Holland Pub. Cie.
\item H. Hironaka, Resolution of singularities of an algebraic variety over a field of characteristic zero, Ann. of Math. 79, pp. 109--326 (1964).
\item L. Hodgkin, tesis (¿de próxima aparición?).
\item L. Illusie, trabajo en preparación.
\item J.-P. Serre, \emph{Algèbre locale. Multiplicités}, Lecture Notes in Mathematics n.º 11, Springer, 1965.
\item A. Borel, Sous-groupes commutatifs et torsion des groupes de Lie compacts connexes, Tohoku Math. Journal, vol. 13, pp. 216--240 (1961).
\item[{[16]}(*)] P. Deligne, Exp. XVIII en \emph{Théorie des topos et cohomologie étale des schémas (SGA \sisi{6}{4})}, por M. Artin, A. Grothendieck y J. L. Verdier. % Canon SGA-SRC-001457
\item[{[17]}(*)] \sisi{U. Mausin}{Yu. I. Manin}, Correspondances, motifs et transformations monoïdales (en ruso). Mat. Sbornik, t. 77 (119), n\textsuperscript{o} 4, 1968. % Canon SGA-SRC-001458
\end{enumerate}

\noindent{\footnotesize (*) Añadido en febrero de 1969.}

\subsection*{Nota para la página 1}

La teoría de las operaciones $\Lambda^i$ en $D^-(X)$ y de las operaciones
$\lambda^i$ en $K'(X)$ fue escrita por P. Deligne en 1968. Demuestra que
$K'(X)$ se convierte así en un pre-$\lambda$-anillo y, cuando
$\mathcal O_X$ es un haz de álgebras sobre un cuerpo, que $K'(X)$ es incluso
un $\lambda$-anillo. Es plausible que un refinamiento del método de Deligne
permita prescindir de la hipótesis restrictiva sobre $\mathcal O_X$.

\subsection*{Notas para la página 11}

\noindent(*) (Abril de 1969) D. MUMFORD nos señala una demostración de (9.6)
en el anillo de Chow (sin despreciar la torsión), que obtuvo en un seminario
en Harvard en 1959. El lector encontrará la demostración de este bello
resultado en SGA 5 VII 9. Este resultado contiene evidentemente el resultado
cohomológico correspondiente, señalado en la nota (**). Sin embargo, el
resultado cohomológico es demostrado por J. P. JOUANOLOU bajo condiciones
más generales (\sisi{sin hiperplano cuasiproyectivo}{sin hipótesis de cuasiproyectividad} y en el caso liso relativo % Canon SGA-SRC-001459
sobre una base cualquiera).

\medskip

\noindent(**) (Enero de 1969) J. P. JOUANOLOU acaba de demostrar estas
fórmulas para todo $y$ sin despreciar la torsión; cf. SGA 5 VII \S{} 4.9.
DELIGNE y MANIN \sisi{17}{[17]} habían observado anteriormente, de manera independiente y % Canon SGA-SRC-001460
mediante argumentos (fáciles) de naturaleza bastante diferente, que bastaba
con establecer fórmulas para $y$ «~algebraico~», e incluso para $y=1$ en \sisi{16}{[16]},
para concluir la misma relación en el caso general.



% Final-index helper macros
\providecommand{\Qcoh}{\operatorname{Qcoh}}
\providecommand{\Coh}{\operatorname{Coh}}
\providecommand{\Parf}{\operatorname{Parf}}
\providecommand{\Lib}{\operatorname{Lib}}
\providecommand{\Loclib}{\operatorname{Loclib}}
\providecommand{\Proj}{\operatorname{Proj}}
\providecommand{\Mod}{\operatorname{Mod}}
\providecommand{\Ob}{\operatorname{Ob}}
\providecommand{\Hom}{\operatorname{Hom}}
\providecommand{\toramp}{\operatorname{tor.amp}}
\providecommand{\parfamp}{\operatorname{par.amp}}
\providecommand{\coh}{\operatorname{coh}}
\providecommand{\qcoh}{\operatorname{qcoh}}
\providecommand{\parf}{\operatorname{parf}}
\providecommand{\torf}{\operatorname{torf}}
\providecommand{\amp}{\operatorname{amp}}
\providecommand{\num}{\operatorname{num}}
\providecommand{\naif}{\operatorname{naif}}
\providecommand{\alg}{\operatorname{alg}}
\providecommand{\Pic}{\operatorname{Pic}}
\providecommand{\Gr}{\operatorname{Gr}}
\providecommand{\Fil}{\operatorname{Fil}}
\providecommand{\Ch}{\operatorname{Ch}}
\providecommand{\Todd}{\operatorname{Todd}}
\providecommand{\Diff}{\operatorname{Diff}}
\providecommand{\Dcat}{\mathrm D}
\providecommand{\Tr}{\operatorname{Tr}}
\providecommand{\Z}{\mathbb Z}
\providecommand{\Q}{\mathbb Q}

\newpage
